\documentclass[12pt,a4paper]{article}

\usepackage[T2A]{fontenc}
\usepackage[utf8]{inputenc}
\usepackage[ukrainian]{babel}
\usepackage[left=25mm,right=15mm,top=20mm,bottom=20mm]{geometry}
\usepackage{longtable}
\usepackage{array}
\usepackage{enumitem}
\usepackage{hyperref}
\hypersetup{colorlinks=true,linkcolor=blue,urlcolor=blue}
\usepackage{tabularx}
\usepackage{booktabs}

\renewcommand{\arraystretch}{1.25}
\setlist{nosep,leftmargin=*}

\newcommand{\req}[2]{\item[\textbf{#1}] #2}

\title{\textbf{Специфікація вимог до програмного забезпечення (SRS)}\\[4pt]
\large Інформаційна система обліку витрат на харчування}
\author{Майстренко Ілля Ігорович, група КН-32}
\date{Версія 1.0 (базова) \\ 29 вересня 2026 р.}

\begin{document}
\maketitle

\begin{center}
\begin{tabular}{@{}ll@{}}
\toprule
Статус & Базова версія, затверджена на SPEC-GATE \\
Репозиторій & \url{https://github.com/d1aboli/food-expenses-is} \\
Вхідні артефакти & \texttt{spec/idea.md}, \texttt{logs/sked-session-01.md}, \texttt{logs/conflict.md} \\
Глосарій & \texttt{spec/glossary.md} \\
Простежуваність & \texttt{spec/traceability.md} \\
\bottomrule
\end{tabular}
\end{center}

\tableofcontents
\newpage

\section{Вступ}

\subsection{Призначення документа}
Документ визначає вимоги до інформаційної системи обліку витрат на харчування. Він є еталонним
артефактом SDD-циклу: на його основі виконується декомпозиція задач, створюються тести
і реалізується функціональність. У разі розбіжності між кодом і цим документом еталоном вважається документ.

\subsection{Область застосування}
Система призначена для особистого обліку щоденних витрат на харчування. Вона дозволяє записувати покупки,
переглядати їх список і отримувати підсумки витрат за обраний період з розбивкою за категоріями.

\subsection{Терміни та визначення}
Терміни вживаються у значенні, наведеному в глосарії \texttt{spec/glossary.md}. Ключові:
\textit{профіль}, \textit{покупка}, \textit{ціна за одиницю}, \textit{кількість}, \textit{вартість покупки},
\textit{категорія}, \textit{календарний тиждень}, \textit{календарний місяць}, \textit{підсумок за період}.

\subsection{Джерела вимог}
Вимоги сформовано на основі початкового задуму (\texttt{spec/idea.md}) та SKED-діалогу,
підсумок якого наведено в \texttt{logs/sked-session-01.md}.

\section{Загальний опис}

\subsection{Призначення системи}
Людина витрачає гроші на їжу щодня дрібними сумами і в кінці місяця не може сказати, куди вони пішли.
Система дозволяє фіксувати кожну покупку за кілька секунд і бачити підсумки за день, тиждень, місяць
або довільний період.

\subsection{Категорії користувачів}
Єдина категорія --- людина, яка веде облік власних витрат. На одному комп'ютері може існувати
кілька профілів; кожен профіль працює лише зі своїми даними.

\subsection{Середовище виконання}
Веб-застосунок, який запускається локально. Серверна частина --- Python (FastAPI),
дані зберігаються в локальній базі SQLite, графіки будуються бібліотекою Chart.js у браузері.

\subsection{Припущення та залежності}
\begin{itemize}
  \item Усі суми вказуються в гривнях.
  \item Профілі не захищені паролем: вони розмежовують дані, а не захищають їх від сторонніх осіб.
  \item Дані вводяться вручну; зовнішні джерела даних не використовуються.
\end{itemize}

\section{Межі системи}

\subsection{Входить до поточної версії}
\begin{itemize}
  \item Створення профілів і перемикання між ними.
  \item Додавання, редагування та видалення покупок.
  \item Перегляд списку покупок.
  \item Базові та власні категорії.
  \item Підсумки за календарний день, тиждень, місяць і довільний діапазон дат.
  \item Розбивка витрат за категоріями.
  \item Стовпчикові графіки по днях місяця та по місяцях року.
  \item Експорт покупок у CSV і резервна копія бази даних.
\end{itemize}

\subsection{Не входить до поточної версії}
\begin{itemize}
  \item Автентифікація за паролем, ролі та права доступу.
  \item Підключення до банківських рахунків, імпорт виписок, сканування чеків.
  \item Планування бюджету, ліміти витрат і сповіщення про перевищення.
  \item Прогнозування витрат і рекомендації.
  \item Синхронізація між пристроями та робота через інтернет.
  \item Кілька валют і курси валют.
\end{itemize}

\subsection{Зміни меж відносно початкового задуму}
Під час SKED-діалогу межі системи розширено трьома рішеннями людини: багатопрофільність,
проста аналітика (розбивка за категоріями і графіки) та облік кількості товару з ціною за одиницю.
Зміни внесено до \texttt{spec/idea.md}.

\section{Функціональні вимоги}

\subsection{Профілі}
\begin{description}
\req{REQ-F-001}{Система дозволяє створити профіль, указавши його назву.}
\req{REQ-F-002}{Система відображає список наявних профілів і дозволяє обрати активний профіль.}
\req{REQ-F-003}{Система показує та змінює лише дані активного профілю: покупки і категорії інших профілів недоступні.}
\req{REQ-F-004}{Якщо жодного профілю не існує, система пропонує створити перший профіль перед роботою з покупками.}
\end{description}

\subsection{Покупки}
\begin{description}
\req{REQ-F-005}{Система дозволяє додати покупку, зазначивши назву, категорію, ціну за одиницю та кількість.}
\req{REQ-F-006}{Система обчислює вартість покупки як добуток ціни за одиницю та кількості, заокруглений до двох знаків після коми за правилом математичного заокруглення (половина заокруглюється вгору).}
\req{REQ-F-007}{Під час додавання покупки система автоматично проставляє поточну дату й час.}
\req{REQ-F-008}{Система дозволяє змінити будь-яке поле раніше доданої покупки, зокрема дату й час.}
\req{REQ-F-009}{Після зміни або видалення покупки система перераховує підсумки за всі періоди, яких ця покупка стосувалася.}
\req{REQ-F-010}{Система дозволяє видалити покупку.}
\req{REQ-F-011}{Система дозволяє додавати покупки з однаковою назвою, категорією та датою; такі записи вважаються різними покупками.}
\req{REQ-F-012}{Система відображає список покупок активного профілю, упорядкований за датою й часом від новіших до старіших.}
\end{description}

\subsection{Перевірка даних}
\begin{description}
\req{REQ-F-013}{Система відхиляє покупку, у якій ціна за одиницю менша або дорівнює нулю.}
\req{REQ-F-014}{Система відхиляє покупку, у якій кількість менша або дорівнює нулю.}
\req{REQ-F-015}{Система відхиляє покупку з порожньою назвою.}
\req{REQ-F-016}{Система відхиляє покупку, дата й час якої пізніші за поточний момент.}
\req{REQ-F-017}{Система приймає ціну як ціле число або число з не більш ніж двома знаками після коми; верхня межа ціни не встановлюється.}
\req{REQ-F-018}{Система приймає кількість як ціле число або число з не більш ніж трьома знаками після коми (наприклад, 0,7 кг або 0,125 кг).}
\req{REQ-F-019}{У разі відхилення система показує повідомлення з причиною і не створює запис.}
\end{description}

\subsection{Категорії}
\begin{description}
\req{REQ-F-020}{При створенні профілю система створює базові категорії: продукти, ресторан, доставка, напої, інше.}
\req{REQ-F-021}{Система дозволяє створити власну категорію з унікальною в межах профілю назвою.}
\req{REQ-F-022}{Система дозволяє перейменувати категорію; покупки зберігають зв'язок з нею.}
\req{REQ-F-023}{Система забороняє видалення непорожньої категорії і показує попередження з кількістю пов'язаних покупок.}
\req{REQ-F-024}{Якщо користувач підтверджує видалення непорожньої категорії, система переносить її покупки до службової категорії <<Інше>> і лише потім видаляє категорію.}
\req{REQ-F-025}{Категорія <<Інше>> є службовою: її не можна ані видалити, ані перейменувати, оскільки до неї переносяться покупки з видалених категорій.}
\end{description}

\subsection{Підсумки та аналітика}
\begin{description}
\req{REQ-F-026}{Система обчислює підсумок витрат за календарний день.}
\req{REQ-F-027}{Система обчислює підсумок витрат за календарний тиждень з понеділка по неділю.}
\req{REQ-F-028}{Система обчислює підсумок витрат за календарний місяць з 1-го по останнє число.}
\req{REQ-F-029}{Система обчислює підсумок витрат за довільний діапазон дат, включно з обома межами.}
\req{REQ-F-030}{Система відображає розбивку витрат за категоріями для обраного періоду у вигляді таблиці <<категорія --- сума --- частка у відсотках>>.}
\req{REQ-F-031}{Система відображає стовпчиковий графік витрат за днями обраного місяця: один стовпчик на кожен день місяця, висота стовпчика дорівнює підсумку за цей день.}
\req{REQ-F-032}{Система відображає стовпчиковий графік витрат за місяцями обраного року: один стовпчик на кожен місяць, висота стовпчика дорівнює підсумку за цей місяць.}
\req{REQ-F-033}{Якщо за обраний період покупок немає, система повертає підсумок 0,00 грн і показує порожній графік.}
\end{description}

\subsection{Зберігання та експорт даних}
\begin{description}
\req{REQ-F-034}{Система зберігає профілі, категорії та покупки в локальній базі даних; дані доступні після перезапуску.}
\req{REQ-F-035}{Система дозволяє експортувати покупки активного профілю за обраний період у файл CSV.}
\req{REQ-F-036}{Система дозволяє створити резервну копію файлу бази даних.}
\end{description}

\section{Нефункціональні вимоги}

\subsection{Процес розробки: TDD}
\begin{description}
\req{REQ-NF-001}{Розробка ведеться за TDD: для кожної функціональної вимоги спочатку створюється автоматизований тест, який не проходить, і лише після цього пишеться код, який його проходить.}
\req{REQ-NF-002}{Кожна функціональна вимога REQ-F-XXX покрита щонайменше одним автоматизованим тестом; ідентифікатор вимоги зазначається в назві або описі тесту.}
\req{REQ-NF-003}{Повний набір тестів запускається однією командою і завершується без помилок перед кожним комітом з префіксом \texttt{feat:}.}
\req{REQ-NF-004}{В історії Git коміт з префіксом \texttt{test:} для вимоги передує коміту з префіксом \texttt{feat:} для тієї самої вимоги.}
\end{description}

\subsection{Процес розробки: журналювання}
\begin{description}
\req{REQ-NF-005}{Результат кожного прогону тестів записується у файл журналу в каталозі \texttt{/logs} із зазначенням дати й часу, загальної кількості тестів, кількості успішних і неуспішних.}
\req{REQ-NF-006}{Кожна сесія роботи з кодинговим агентом фіксується у \texttt{/logs}: запит людини, перелік змінених файлів, результат перегляду через \texttt{git diff} і рішення про прийняття або відхилення змін.}
\req{REQ-NF-007}{Зміни, згенеровані кодинговим агентом, не додаються до репозиторію без перегляду людиною; повідомлення відповідного коміту містить позначку \texttt{[agent]}.}
\req{REQ-NF-008}{Помилки виконання застосунку записуються у файл журналу в каталозі \texttt{/logs} із датою, часом і текстом помилки.}
\end{description}

\subsection{Якість роботи системи}
\begin{description}
\req{REQ-NF-009}{Підсумок за будь-який період обчислюється не довше ніж за 1 секунду для бази, що містить до 10\,000 покупок.}
\req{REQ-NF-010}{Грошові суми зберігаються й відображаються з точністю до двох знаків після коми; проміжні обчислення не втрачають копійки.}
\req{REQ-NF-011}{Покупка записується до бази даних у момент підтвердження форми, тому аварійне завершення застосунку не призводить до втрати вже збережених покупок.}
\req{REQ-NF-012}{Усі тексти інтерфейсу та повідомлення про помилки зберігаються у файлах перекладів для української та англійської мов; у коді немає текстів, зашитих безпосередньо.}
\end{description}

\subsection{Обмеження реалізації}
\begin{description}
\req{REQ-NF-013}{Серверна частина реалізується мовою Python із використанням FastAPI; дані зберігаються в SQLite.}
\req{REQ-NF-014}{Графіки будуються бібліотекою Chart.js на боці браузера.}
\req{REQ-NF-015}{Код розміщується в каталозі \texttt{/src}, тести --- у каталозі \texttt{/tests} відповідно до структури репозиторію, визначеної в лабораторній роботі №1.}
\end{description}

\section{Вимоги до інтерфейсів}
\begin{description}
\req{REQ-I-001}{Користувацький інтерфейс --- вебсторінка з формою додавання покупки, списком покупок, вибором періоду, таблицею розбивки за категоріями та графіком.}
\req{REQ-I-002}{Вибір активного профілю доступний з будь-якої сторінки застосунку.}
\req{REQ-I-003}{Формат експорту CSV: назва, категорія, ціна за одиницю, кількість, вартість, дата й час; роздільник --- кома, кодування UTF-8.}
\req{REQ-I-004}{Дати відображаються у форматі ДД.ММ.РРРР, час --- у форматі ГГ:ХХ.}
\end{description}

\section{Обмеження}
\begin{itemize}
  \item Єдина валюта --- гривня.
  \item Робота лише на локальному комп'ютері, без доступу через інтернет.
  \item Профілі без пароля; система не призначена для зберігання конфіденційних даних.
  \item Обсяг даних одного профілю --- до 10\,000 покупок.
\end{itemize}

\section{Критерії прийняття}

\begin{longtable}{@{}p{0.17\textwidth}p{0.77\textwidth}@{}}
\toprule
\textbf{Вимога} & \textbf{Критерій прийняття} \\
\midrule
\endhead
REQ-F-003 & У профілі <<Ілля>> створено покупку; після перемикання на профіль <<Гість>> список покупок порожній. \\
REQ-F-006 & Для ціни 80,00 грн за кілограм і кількості 0,7 вартість покупки дорівнює 56,00 грн. \\
REQ-F-007 & Після додавання покупки її дата й час збігаються з системним часом на момент створення з точністю до хвилини. \\
REQ-F-009 & Після зміни дати покупки зі 05.10 на 04.10 підсумок за 05.10 зменшується на її вартість, а підсумок за 04.10 зростає на ту саму суму. \\
REQ-F-013 & При спробі додати покупку з ціною 0 або -50 запис не створюється, показується повідомлення про помилку, кількість покупок не змінюється. \\
REQ-F-016 & При спробі встановити дату й час, пізніші за поточний момент, запис не зберігається. \\
REQ-F-018 & Кількість 0,125 приймається; кількість 0,1255 відхиляється. \\
REQ-F-023 & При спробі видалити категорію, до якої належать 3 покупки, показується попередження із зазначенням кількості покупок. \\
REQ-F-024 & Після підтвердження видалення непорожньої категорії всі її покупки мають категорію <<Інше>>, а сама категорія відсутня у списку. \\
REQ-F-025 & Кнопки видалення та перейменування для категорії <<Інше>> недоступні. \\
REQ-F-027 & Для покупок 100,00 грн у понеділок і 50,00 грн у неділю того самого тижня підсумок за тиждень дорівнює 150,00 грн; покупка наступного понеділка в підсумок не входить. \\
REQ-F-029 & Для діапазону 01.10--03.10 у підсумок входять покупки за 01.10 і 03.10 включно. \\
REQ-F-030 & Сума часток за категоріями для непорожнього періоду дорівнює 100\% з точністю до 0,1\%. \\
REQ-F-031 & Для жовтня 2026 графік містить 31 стовпчик; висота стовпчика за 04.10 дорівнює підсумку за 04.10. \\
REQ-F-032 & Для 2026 року графік містить 12 стовпчиків; висота стовпчика за жовтень дорівнює підсумку за жовтень. \\
REQ-F-033 & Для місяця без покупок підсумок дорівнює 0,00 грн, графік не містить стовпчиків. \\
REQ-F-034 & Після перезапуску застосунку раніше додані покупки присутні у списку. \\
REQ-F-035 & Експортований файл CSV містить рядок заголовків і по одному рядку на кожну покупку періоду; суми у файлі збігаються зі списком. \\
REQ-NF-001 & Для кожної реалізованої вимоги в історії Git наявний коміт \texttt{test:}, датований раніше за відповідний коміт \texttt{feat:}. \\
REQ-NF-002 & Запуск набору тестів показує щонайменше один тест для кожного ідентифікатора REQ-F-XXX. \\
REQ-NF-005 & Після прогону тестів у каталозі \texttt{/logs} існує файл із датою прогону і рядком виду <<усього N, успішних M, неуспішних K>>. \\
REQ-NF-009 & На базі з 10\,000 покупок обчислення підсумку за місяць триває менше ніж 1 секунду. \\
REQ-NF-010 & Для ціни 33,33 грн і кількості 3 вартість дорівнює 99,99 грн, а не 100,00 грн. \\
REQ-NF-012 & Кожен текстовий рядок інтерфейсу присутній у файлах перекладів для обох мов; пошук по коду не знаходить текстів поза цими файлами. \\
\bottomrule
\end{longtable}

\section{Сценарії поведінки (BDD)}
Повний перелік сценаріїв наведено у файлі \texttt{spec/scenarios/expenses.feature}. Нижче --- ключові.

\subsection*{Сценарій 1. Додавання покупки з кількістю}
\begin{itemize}
  \item \textbf{Given} обрано профіль <<Ілля>>, список покупок порожній
  \item \textbf{When} користувач додає покупку <<Яблука>>, категорія <<Продукти>>, ціна 80,00 грн за кілограм, кількість 0,7
  \item \textbf{Then} у списку одна покупка з вартістю 56,00 грн
  \item \textbf{And} її дата й час дорівнюють поточним
\end{itemize}

\subsection*{Сценарій 2. Відмова від'ємної ціни}
\begin{itemize}
  \item \textbf{Given} список покупок порожній
  \item \textbf{When} користувач додає покупку з ціною -50,00 грн
  \item \textbf{Then} покупка не створюється
  \item \textbf{And} показується повідомлення <<Ціна має бути більшою за нуль>>
\end{itemize}

\subsection*{Сценарій 3. Межа календарного тижня}
\begin{itemize}
  \item \textbf{Given} є покупка 100,00 грн за неділю 04.10.2026
  \item \textbf{And} є покупка 40,00 грн за понеділок 05.10.2026
  \item \textbf{When} користувач переглядає підсумок за тиждень, що містить 04.10.2026
  \item \textbf{Then} підсумок дорівнює 100,00 грн
\end{itemize}

\subsection*{Сценарій 4. Видалення непорожньої категорії}
\begin{itemize}
  \item \textbf{Given} у категорії <<Напої>> є 3 покупки
  \item \textbf{When} користувач видаляє категорію <<Напої>>
  \item \textbf{Then} показується попередження про 3 пов'язані покупки
  \item \textbf{And} після підтвердження ці покупки мають категорію <<Інше>>
\end{itemize}

\subsection*{Сценарій 5. Ізоляція профілів}
\begin{itemize}
  \item \textbf{Given} у профілі <<Ілля>> є 2 покупки
  \item \textbf{When} користувач перемикається на профіль <<Гість>>
  \item \textbf{Then} список покупок порожній
  \item \textbf{And} підсумок за поточний місяць дорівнює 0,00 грн
\end{itemize}

\subsection*{Сценарій 6. Перерахунок після зміни дати}
\begin{itemize}
  \item \textbf{Given} є покупка 60,00 грн за 05.10.2026
  \item \textbf{And} підсумок за 05.10.2026 дорівнює 60,00 грн
  \item \textbf{When} користувач змінює дату покупки на 04.10.2026
  \item \textbf{Then} підсумок за 05.10.2026 дорівнює 0,00 грн
  \item \textbf{And} підсумок за 04.10.2026 дорівнює 60,00 грн
\end{itemize}

\section{Простежуваність}
Початкова матриця простежуваності зберігається у файлі \texttt{spec/traceability.md}.
Вона пов'язує ідентифікатор вимоги з критерієм прийняття та сценарієм; колонки для тікетів,
тестів, коду й документації створено заздалегідь і заповнюватимуться в наступних лабораторних роботах.

\section{Відкриті питання}
\begin{description}
\req{ВП-01}{Чи потрібне видалення профілю разом з усіма його покупками і категоріями? Питання виникло під час SKED-діалогу і рішення не прийнято. До його закриття система видалення профілю не виконує.}
\end{description}

\section{Історія версій}
\begin{tabular}{@{}lll@{}}
\toprule
Версія & Дата & Зміни \\
\midrule
1.0 & 29.09.2026 & Базова специфікація, затверджена на SPEC-GATE \\
\bottomrule
\end{tabular}

\end{document}
